\section{Síntesis y ampliación}

Esta clase conectó de principio a fin tres líneas de trabajo: usar un learned model para ampliar los datos, usar el condicionamiento por tarea para construir una generalist policy y usar hindsight relabeling para convertir las trayectorias fallidas en supervisión útil. Aunque en apariencia son distintas, en esencia las tres responden a la misma pregunta: \textbf{cómo hacer que una experiencia limitada sirva para más decisiones.}

\subsection{Tabla resumen de la clase}

\begin{table}[H]
\centering
\begin{tabular}{p{0.18\textwidth} p{0.22\textwidth} p{0.22\textwidth} p{0.24\textwidth}}
\toprule
Tema & Pregunta central & Técnica principal & Limitación principal \\
\midrule
Model-based RL & Cómo reducir el costo de la interacción real & learned model + short synthetic rollout & Sesgo del modelo, acumulación de error en horizons largos \\
RL multitarea & Cómo compartir experiencia y habilidades entre tareas & task conditioning + shared policy & Transferencia negativa, distribución desbalanceada de tareas \\
Goal-conditioned RL & Cómo aprender una política general para alcanzar metas & La meta entra en la política como condición & Dificultad para representar la meta y diseñar la recompensa \\
HER & Cómo lograr que las trayectorias fallidas también sirvan para entrenar & hindsight relabeling + off-policy reuse & Requiere recompensas recalculables y dinámica compartida \\
\bottomrule
\end{tabular}
\caption{Hilo de los métodos de la Lecture 12}
\end{table}

\subsection{Reflexiones adicionales}

\begin{itemize}
    \item Cuando el task descriptor pasa de one-hot a lenguaje natural, la frontera entre el RL multitarea y los foundation models se vuelve cada vez más difusa.
    \item El éxito de HER nos recuerda que muchos datos que parecen «fallidos» en realidad solo carecen de una etiqueta y una interpretación adecuadas.
    \item La pregunta que el curso anticipa al final surge de forma natural: si es posible compartir conocimiento entre varias tareas, ¿se puede ir más allá y lograr una adaptación rápida a tareas completamente nuevas?
\end{itemize}

\subsection{Del aprendizaje compartido entre tareas a la adaptación rápida}

El final de esta clase deja planteado un giro de investigación muy importante: si el RL multitarea resuelve el problema de «compartir estructura entre un conjunto de tareas conocidas», el paso siguiente es preguntarse: \textbf{ante una tarea nueva que nunca se ha visto, ¿es posible adaptarse rápidamente, como en el in-context learning, en lugar de volver a entrenar toda la política?}

\begin{figure}[H]
\centering
\includegraphics[width=0.9\textwidth]{slides-images/slide-29.jpg}
\caption{La siguiente pregunta que plantea el cierre de la clase: ¿es posible lograr una adaptación rápida a tareas nuevas?}
\end{figure}

Desde esta perspectiva, las tres líneas principales de esta clase constituyen precisamente la base de la investigación posterior:
\begin{enumerate}
    \item El Model-based RL permite que el agent reutilice mejor la experiencia existente y la estructura del entorno.
    \item El RL multitarea permite que la política aprenda a compartir parámetros y patrones de comportamiento entre varias tareas.
    \item HER muestra que una misma trayectoria puede aportar supervisión repetidamente bajo distintas interpretaciones de la meta.
\end{enumerate}

En conjunto apuntan a una pregunta más amplia: \textit{si la experiencia, el condicionamiento por tarea y la interpretación de la meta pueden reutilizarse, ¿no debería la política contar también con una mayor capacidad de test-time adaptation?}

\subsection{Errores más frecuentes al implementar los métodos de esta clase}

Aunque conceptualmente esta clase se divide en tres partes, al escribir el código los errores suelen aparecer en el pipeline de datos. En particular, con HER y con las políticas condicionadas por tarea, muchas veces el problema no está en las fórmulas del algoritmo, sino en que el replay buffer, el task descriptor y la reward recomputation no están bien conectados entre sí.

\begin{figure}[H]
\centering
\includegraphics[width=0.9\textwidth]{slides-images/slide-28.jpg}
\caption{Repaso de las líneas principales de la Lecture 12: model-based augmentation, task conditioning y hindsight relabeling}
\end{figure}

\begin{table}[H]
\centering
\begin{tabular}{p{0.2\textwidth} p{0.28\textwidth} p{0.28\textwidth}}
\toprule
Módulo de implementación & Error más común & Práctica más segura \\
\midrule
Synthetic rollout & Un rollout demasiado largo hace que el sesgo del modelo domine el entrenamiento & Partir de estados intermedios de los datos reales y hacer solo rollouts cortos \\
Task conditioning & El task ID y el estado no se envían juntos a la policy / critic & Usar de forma uniforme $(\bar{s}, z)$ como interfaz de entrada condicional \\
HER relabeling & Cambiar solo la meta sin recalcular la recompensa de toda la trayectoria & Implementar explícitamente la reward recomputation y la lógica de reescritura en el buffer \\
Replay buffer & Desbalance entre la distribución de muestras originales y la de muestras relabeled & Registrar la proporción de muestreo y, si hace falta, usar muestreo estratificado \\
\bottomrule
\end{tabular}
\caption{Lista de verificación de implementación de la Lecture 12}
\end{table}

\begin{knowledgebox}{Lecturas adicionales}
\begin{itemize}
    \item Andrychowicz et al., ``Hindsight Experience Replay,'' NeurIPS 2017
    \item Schaul et al., ``Universal Value Function Approximators,'' ICML 2015
    \item Yu et al., ``Meta-World: A Benchmark and Evaluation for Multi-Task and Meta Reinforcement Learning,'' CoRL 2020
    \item Ghosh et al., ``Learning to Reach Goals via Iterated Supervised Learning,'' ICLR 2021
\end{itemize}
\end{knowledgebox}

\end{document}
